% Generated draft: requires mathematical and source review.
\documentclass[11pt]{article}
\usepackage[margin=25mm]{geometry}
\usepackage{fontspec}
\setmainfont{Latin Modern Roman}
\usepackage{amsmath,amssymb,mathtools,unicode-math}
\setmathfont{Latin Modern Math}
\usepackage{enumitem,needspace,xurl,hyperref}
\setlength{\parindent}{0pt}
\setlength{\parskip}{5pt}
\setlength{\emergencystretch}{3em}
\newcommand{\N}{\mathbb{N}}
\newcommand{\Z}{\mathbb{Z}}
\newcommand{\Q}{\mathbb{Q}}
\newcommand{\R}{\mathbb{R}}
\newcommand{\C}{\mathbb{C}}
\newcommand{\F}{\mathbb{F}}
\newcommand{\A}{\mathbb{A}}
\newcommand{\PP}{\mathbb{P}}
\newcommand{\E}{\mathbb{E}}
\newcommand{\eps}{\varepsilon}
\newcommand{\dd}{\,\mathrm{d}}
\begin{document}
\section*{UnsolvedMath ID: 30001733}
\textbf{Source problem code:} catalogue entry
\subsection{Definitions and mathematical statement}
Let \(T\) denote a tree used in the source's smooth tree-based level-set estimation procedure, and let \(\rho\) denote the parameter controlling the penalty for tree complexity. These symbols specify their roles only: the supplied record does not define the estimator, the tree class, or the penalty.

\textbf{Supplied qualitative conjecture (2011).} In smooth tree-based level-set estimation, if relatively simple trees with few branches suffice to obtain a good estimate, then the choice of the tree-complexity penalty parameter \(\rho\) is not crucial for obtaining a good estimate.

\textbf{Unresolved mathematical specification.} The supplied record does not specify the underlying domain, the function or distribution whose level set is estimated, the level or boundary convention, the observation model, or the meaning of smoothness. It also leaves undefined the relevant measure of tree complexity, the threshold for having few branches, the estimation loss, and the criterion for a good estimate. The admissible values of \(\rho\), the quantifiers governing its choice, and the meaning of its being inessential are unspecified. In particular, the record does not determine whether the proposed robustness concerns identical estimates, comparable finite-sample errors, or an asymptotic property. Consequently, this qualitative conjecture cannot be converted into a uniquely specified quantified mathematical proposition from the supplied evidence.
\subsection{Short English statement}
The conjecture asks whether the tree-complexity penalty parameter \(\rho\) has little effect on estimation quality when a tree with few branches already gives a good level-set estimate. The estimator and the meanings of simplicity, accuracy, and robustness are unspecified, so a precise quantitative statement requires further source information.
\subsection{Sources}
\begin{itemize}
\item Local UnsolvedMath registry: supplied complete source record for permanent ID \(30001733\), titled ``Penalty Robustness in Tree-Based Level-Set Estimation,'' source problem code \texttt{catalogue entry}, in the Oberwolfach Reports Open Problems collection. The supplied proposed year is 2011; no individual proposer is identified.
\item \emph{Mini-Workshop: Level Sets and Depth Contours in High Dimensional Data} (2011), Oberwolfach Reports, pp.\ 691--717, DOI \texttt{10.4171/owr/2011/13}. Supplied source and literature reference, labeled ``Oberwolfach Report 13/2011'': \url{https://doi.org/10.4171/owr/2011/13}. The record describes Oberwolfach Reports as documenting mathematical workshops and their open-problem discussions; no precise page within the report is supplied for this conjecture.
\item Supplied dated literature curation, checked \texttt{2026-08-21}: the status is labeled open, and the assessment states that the report presents the penalty-irrelevance claim for simple trees as a conjecture and that no general theorem proving or disproving it was located. The listed source checked is the report linked above. This is supplied dated curation, not an independently verified literature assessment; time-sensitive status and cited claims require independent verification.
\item The registry labels the statement \texttt{corrected\_verified} and reports that its mathematical content was checked against the official report, with notation, context, and question form normalized without changing the source-level claim. This verification description is supplied provenance and has not been independently checked for this draft.
\end{itemize}
\end{document}
